\documentclass[../../main.tex]{subfiles}
\begin{document}

{\bf [解答]}

与えられた数列
\begin{align}
  c_n = (n+1) \int_0^1 x^n \cos\pi x \, dx \label{eq:1}
\end{align}
について考える．

\paragraph{(1)}

\cref{eq:1}に$\cos\pi x$に対する部分積分を用いて$c_{n+2}$を変形すると
\begin{align}
  c_{n+2}
  &= (n+3) \int_0^1 x^{n+2} \cos\pi x \, dx \\
  &= (n+3) \left[ \left[ \dfrac{1}{\pi} \sin\pi x \cdot x^{n+2} \right]_0^1 - \int_0^1 \dfrac{n+2}{\pi} x^{n+1} \sin\pi x \, dx \right] \\
  &= - \dfrac{1}{\pi} (n+3)(n+2) \int_0^1 x^{n+1} \sin\pi x \, dx \label{eq:2} \\
  &= - \dfrac{1}{\pi} (n+3)(n+2) \left[ \left[ \dfrac{-1}{\pi} \cos\pi x \cdot x^{n+1} \right]_0^1 + \dfrac{1}{\pi} c_n \right] \\
  &= - \dfrac{(n+3)(n+2)}{\pi} \left\{ \dfrac{1}{\pi} + \dfrac{1}{\pi} c_n \right\} \\
  &= - \dfrac{(n+3)(n+2)}{\pi^2} \left\{ 1 + c_n \right\}
\end{align}

\paragraph{(2)}

\cref{eq:1}を今度は$x^{n}$について部分積分して
\begin{align}
  c_n
  &= (n+1)\int_0^1 x^n \cos\pi x \, dx \\
  &= \left[ x^{n+1} \cos\pi x \right]_0^1 - \int_0^1 x^{n+1} (-\pi \sin\pi x) \, dx \\
  &= -1 + \pi \int_0^1 x^{n+1} \sin\pi x \, dx \label{eq:3}
\end{align}
である．ここで右辺第二項について，$[0,1]$において$0 \le \sin\pi x \le 1$だから
\begin{align}
  0 \le x^{n+1} \sin\pi x \le x^{n+1}
\end{align}
なので，これを同区間$[0,1]$で積分すると
\begin{align}
  0 \le \int_0^1 x^{n+1} \sin\pi x \, dx \le \int_0^1 x^{n+1} \, dx = \dfrac{1}{n+2}
\end{align}
だから，\cref{eq:3}より
\begin{align}
  0 \le c_n + 1 \le \dfrac{\pi}{n+2} \label{eq:4}
\end{align}
である．右辺は$n\to\infty$で$0$に収束するから，挟み撃ちの定理より
\begin{align}
  \lim_{n\to\infty} c_n = -1 \label{eq:5}
\end{align}
である．

\paragraph{(3)}

\cref{eq:5}より$c=-1$である．従って\cref{eq:3}より
\begin{align}
  c_n - c = \pi \int_0^1 x^{n+1} \sin\pi x \, dx
\end{align}
である．従って$a_n = \dfrac{c_{n+1}-c}{c_n-c}$ とおくと，
\begin{align}
  a_n = \dfrac{\int_0^1 x^{n+2} \sin\pi x \, dx}{\int_0^1 x^{n+1} \sin\pi x \, dx} \label{eq:6}
\end{align}
となる．ここで\cref{eq:2}より
\begin{align}
  \int_0^1 x^{n+1} \sin\pi x \, dx = -\dfrac{\pi c_{n+2}}{(n+3)(n+2)}
\end{align}
だから，\cref{eq:6}に代入して求める極限値は
\begin{align}
  a_n
  & = \dfrac{\pi c_{n+3}/(n+4)(n+3)}{\pi c_{n+2}/(n+3)(n+2)} \\
  &= \dfrac{n+2}{n+4} \cdot \dfrac{c_{n+3}}{c_{n+2}} \\
  &= \dfrac{1+2/n}{1+4/n} \cdot \dfrac{c_{n+3}}{c_{n+2}} \xrightarrow{n\to\infty} 1 \quad \ (\because c_n \to -1)
\end{align}
である．

\end{document}
